\section{Conjuntos numerables}

\begin{definition}[Conjunto numerable]
  Un conjunto \(X\) es \emph{numerable} si \(X \sim \N\). Cada biyección \(\N \to X\) se llama una \emph{enumeración} de \(X\). Un conjunto es \emph{a lo sumo numerable} si es finito o numerable.
\end{definition}

\begin{eg}
  El conjunto de los números pares \(P = \{ 2n : n \in \N \}\) es numerable, pues \(f\colon \N \to P\), \(f(n) = 2n\), es biyectiva.
\end{eg}

\begin{eg}
  \(\Z\) es numerable. La función \(f\colon \Z \to \N\) dada por
  \[
    f(n) =
    \begin{cases}
      2n      & \text{si } n > 0,    \\
      -2n + 1 & \text{si } n \leq 0,
    \end{cases}
  \]
  es biyectiva (los positivos van a los pares y los no positivos a los impares), y su inversa es una enumeración de \(\Z\).
\end{eg}

\begin{theorem}[Cantor]
  Sea \(X\) un conjunto y \(\mathcal{P}(X) = \{ A : A \subset X \}\). Entonces \(\card(X) \neq \card(\mathcal{P}(X))\).
\end{theorem}

\begin{proof}
  Supongamos que existe \(f\colon X \to \mathcal{P}(X)\) biyectiva. Para cada \(x \in X\) puede ocurrir \(x \in f(x)\) o \(x \notin f(x)\). Sea
  \[
    B = \{ x \in X : x \notin f(x) \} \subset X.
  \]
  Como \(f\) es suryectiva, existe \(y \in X\) tal que \(f(y) = B\).
  \begin{itemize}
    \item Si \(y \in B = f(y)\), entonces \(y \in f(y)\), lo que contradice la definición de \(B\).
    \item Si \(y \notin B = f(y)\), entonces \(y \notin f(y)\), y por definición \(y \in B\).
  \end{itemize}
  En ambos casos llegamos a un absurdo, de modo que no existe tal \(f\).
\end{proof}

\begin{definition}[Orden entre cardinales]
  Se escribe \(\card(X) \leq \card(Y)\) si existe una función inyectiva \(X \to Y\). Se escribe \(\card(X) < \card(Y)\) si \(\card(X) \leq \card(Y)\) pero \(X \not\sim Y\).
\end{definition}

\begin{theorem}[Cantor--Schröder--Bernstein]
  Si existen funciones inyectivas \(f\colon X \to Y\) y \(g\colon Y \to X\), entonces existe una función biyectiva \(h\colon X \to Y\).
\end{theorem}

\begin{proof}
  Basta encontrar particiones \(X = X_1 \cup X_2\) e \(Y = Y_1 \cup Y_2\) tales que \(f\colon X_1 \to Y_1\) y \(g\colon Y_2 \to X_2\) sean biyectivas: entonces
  \[
    h(x) =
    \begin{cases}
      f(x)      & \text{si } x \in X_1, \\
      g^{-1}(x) & \text{si } x \in X_2,
    \end{cases}
  \]
  es biyectiva.

  Definimos \(\phi\colon \mathcal{P}(X) \to \mathcal{P}(X)\) por \(\phi(A) = X \setminus g\bigl(Y \setminus f(A)\bigr)\).

  \emph{\(\phi\) es creciente}, es decir, \(A \subset B \Rightarrow \phi(A) \subset \phi(B)\). En efecto,
  \begin{multline*}
    A \subset B \Rightarrow f(A) \subset f(B) \Rightarrow Y \setminus f(B) \subset Y \setminus f(A) \\
    \Rightarrow g\bigl(Y \setminus f(B)\bigr) \subset g\bigl(Y \setminus f(A)\bigr) \Rightarrow \phi(A) \subset \phi(B).
  \end{multline*}

  Sea \(\mathscr{C} = \{ C \subset X : \phi(C) \subset C \}\). Es no vacío pues \(X \in \mathscr{C}\). Sea \(A = \bigcap_{C \in \mathscr{C}} C\). Para cada \(C \in \mathscr{C}\) se tiene \(A \subset C\), luego \(\phi(A) \subset \phi(C) \subset C\); por lo tanto \(\phi(A) \subset A\). Aplicando \(\phi\) nuevamente, \(\phi(\phi(A)) \subset \phi(A)\), es decir, \(\phi(A) \in \mathscr{C}\) y entonces \(A \subset \phi(A)\). Concluimos que \(A = \phi(A)\).

  Sean \(X_1 = A\), \(X_2 = X \setminus A\), \(Y_1 = f(A)\) e \(Y_2 = Y \setminus f(A)\). De \(A = \phi(A)\) se obtiene \(X \setminus A = g\bigl(Y \setminus f(A)\bigr)\), es decir, \(X_2 = g(Y_2)\). Como \(f\) y \(g\) son inyectivas, \(f\colon X_1 \to Y_1\) y \(g\colon Y_2 \to X_2\) son biyectivas.
\end{proof}

\begin{prop}
  La relación \(\leq\) entre cardinales es reflexiva, transitiva y antisimétrica.
\end{prop}

\begin{proof}
  Reflexiva: la identidad es inyectiva. Transitiva: la composición de funciones inyectivas es inyectiva. Antisimétrica: si \(\card(X) \leq \card(Y)\) y \(\card(Y) \leq \card(X)\), por el teorema de Cantor--Schröder--Bernstein \(X \sim Y\).
\end{proof}

\begin{eg}
  \(\N \sim \Z \sim \Q\).
\end{eg}

\begin{proof}
  Ya vimos que \(\N \sim \Z\). La inclusión \(\Z \to \Q\) es inyectiva. Recíprocamente, todo \(q \in \Q\) se escribe de manera única como \(q = a/b\) con \(a \in \Z\), \(b \in \N\) coprimos. Definimos \(F\colon \Q \to \Z\) por \(F(q) = \operatorname{sgn}(a) \cdot 2^{|a|} \cdot 3^{b}\). Por el teorema fundamental de la aritmética, \(F\) es inyectiva. Por el teorema de Cantor--Schröder--Bernstein, \(\Z \sim \Q\).
\end{proof}

\begin{eg}
  \(\N \times \N\) es numerable.
\end{eg}

\begin{proof}
  La función \(n \mapsto (1, n)\) es inyectiva de \(\N\) en \(\N \times \N\), y la función \((n, m) \mapsto 2^n 3^m\) es inyectiva de \(\N \times \N\) en \(\N\) por el teorema fundamental de la aritmética. Por el teorema de Cantor--Schröder--Bernstein, \(\N \times \N \sim \N\).
\end{proof}

\begin{eg}
  El conjunto de sucesiones de ceros y unos, \(A = {\{0, 1\}}^{\N} = \{ {(a_n)}_{n \geq 1} : a_n \in \{0, 1\} \}\), no es numerable.
\end{eg}

\begin{proof}
  Supongamos que es numerable y escribamos \(A = \{ a^1, a^2, \dots, a^j, \dots \}\) con \(a^j = {(a^j_n)}_{n \geq 1}\). Consideremos la sucesión \(b = {(b_n)}_{n \geq 1}\) dada por \(b_n = 1 - a^n_n\), que en el lugar \(n\) tiene lo contrario de lo que tiene la \(n\)-ésima sucesión de la lista. Entonces \(b \in A\) pero \(b \neq a^n\) para todo \(n\) (difieren en el lugar \(n\)), lo que es absurdo.
\end{proof}

Este argumento (\emph{diagonal de Cantor}) puede adaptarse para probar que el intervalo \((0, 1)\) no es numerable. A continuación damos otra demostración, basada en el siguiente teorema.

\begin{theorem}[Principio de encaje de intervalos]
  Sea \(I_1 \supset I_2 \supset I_3 \supset \cdots\) una sucesión decreciente de intervalos cerrados y acotados, \(I_n = [a_n, b_n]\). Entonces
  \[
    \bigcap_{n=1}^{\infty} I_n = [\alpha, \beta], \qquad \alpha = \sup_n a_n, \quad \beta = \inf_n b_n,
  \]
  y en particular la intersección es no vacía.
\end{theorem}

\begin{proof}
  Como \(I_{n+1} \subset I_n\), se tiene \(a_n \leq a_{n+1} \leq b_{n+1} \leq b_n\). Luego, para todos \(m, k \in \N\), tomando \(j = \max\{m, k\}\) resulta \(a_m \leq a_j \leq b_j \leq b_k\). Entonces el conjunto \(A = \{ a_n : n \in \N \}\) está acotado superiormente por cada \(b_k\), y \(B = \{ b_n : n \in \N \}\) está acotado inferiormente por cada \(a_m\). Así existen \(\alpha = \sup A\) y \(\beta = \inf B\), y \(\alpha \leq b_k\) para todo \(k\), de donde \(\alpha \leq \beta\). Se cumple
  \[
    a_n \leq \alpha \leq \beta \leq b_n \quad \text{para todo } n \in \N,
  \]
  y por lo tanto \([\alpha, \beta] \subset I_n\) para todo \(n\), es decir, \([\alpha, \beta] \subset \bigcap_n I_n\).

  Recíprocamente, si \(x < \alpha\), existe \(a_i\) con \(x < a_i\), de modo que \(x \notin I_i\). Análogamente, si \(x > \beta\) existe \(b_i < x\) y \(x \notin I_i\). Luego \(\bigcap_n I_n \subset [\alpha, \beta]\).
\end{proof}

\begin{theorem}
  \(\R\) no es numerable.
\end{theorem}

\begin{proof}
  Supongamos que existe una enumeración \(\N \to \R\), \(n \mapsto x_n\). Construimos recursivamente intervalos cerrados \([0, 1] = I_0 \supset I_1 \supset I_2 \supset \cdots\) tales que \(I_n\) tiene longitud \(3^{-n}\) y \(x_n \notin I_n\). Dado \(I_{n-1}\), lo dividimos en tres intervalos cerrados de igual longitud; como \(x_n\) pertenece a lo sumo a dos de ellos, alguno no lo contiene, y lo llamamos \(I_n\).

  Por el principio de encaje, existe \(x \in \bigcap_{n \geq 1} I_n \subset [0, 1]\). Pero \(x \neq x_n\) para todo \(n\) (pues \(x \in I_n\) y \(x_n \notin I_n\)), en contradicción con que la enumeración recorre todo \(\R\).
\end{proof}

\begin{note}
  Otra idea para ver que \(\R\) no es numerable es la siguiente. Si \(\R = \{x_n : n \in \N\}\), fijado \(0 < \e < 1\) se asigna a cada \(x_n\) un intervalo abierto \(I_n\) centrado en \(x_n\) de longitud \(\e\, 2^{-n}\). Estos intervalos cubrirían \([0, 1]\) con longitud total \(\sum_n \e\, 2^{-n} = \e < 1\), lo cual es imposible, pues un intervalo cerrado no puede cubrirse con intervalos abiertos de longitud total menor que la suya. Este hecho se probará más adelante.
\end{note}

\section{Propiedades de los conjuntos numerables}

\begin{theorem}
  Si \(X\) es numerable e \(Y \subset X\), entonces \(Y\) es a lo sumo numerable.
\end{theorem}

\begin{proof}
  Supongamos que \(Y\) es infinito y escribamos \(X = \{ x_n : n \in \N \}\) (enumeración). Definimos recursivamente
  \[
    n_1 = \min\{ n \in \N : x_n \in Y \}, \qquad n_{k+1} = \min\{ n > n_k : x_n \in Y \},
  \]
  que están bien definidos porque \(Y\) es infinito. La sucesión \((n_k)\) es estrictamente creciente. Sea \(g\colon \N \to Y\) dada por \(g(k) = x_{n_k}\).

  \(g\) es inyectiva pues \(k \neq l \Rightarrow n_k \neq n_l\). Veamos que es suryectiva: sea \(y = x_j \in Y\). Como \((n_k)\) es estrictamente creciente, existe un menor \(k\) con \(n_k \geq j\). Si \(k = 1\), entonces \(n_1 \geq j\) y, como \(x_j \in Y\), \(n_1 \leq j\); luego \(n_1 = j\). Si \(k > 1\), entonces \(n_{k-1} < j \leq n_k\) y como \(x_j \in Y\), la definición de \(n_k\) como mínimo da \(n_k \leq j\); luego \(n_k = j\). En ambos casos \(y = g(k)\).
\end{proof}

\begin{corollary}
  Si \(f\colon X \to Y\) es inyectiva e \(Y\) es a lo sumo numerable, entonces \(X\) es a lo sumo numerable.
\end{corollary}

\begin{proof}
  Como \(f\) es inyectiva, \(X \sim f(X)\), y \(f(X) \subset Y\) es a lo sumo numerable por el teorema anterior.
\end{proof}

\begin{theorem}
  Si \(f\colon X \to Y\) es suryectiva y \(X\) es a lo sumo numerable, entonces \(Y\) es a lo sumo numerable.
\end{theorem}

\begin{proof}
  Como \(f\) es suryectiva, existe \(g\colon Y \to X\) con \(f \circ g = \id_Y\) (inversa a derecha). Entonces \(g\) es inyectiva y, por el corolario, \(Y\) es a lo sumo numerable.
\end{proof}

\begin{theorem}
  Si \((X_n)_{n \in \N}\) es una familia de conjuntos numerables, entonces \(X = \bigcup_{n \in \N} X_n\) es numerable.
\end{theorem}

\begin{proof}
  Para cada \(n\) elegimos una biyección \(f_n\colon \N \to X_n\). Definimos \(F\colon \N \times \N \to X\) por \(F(n, m) = f_n(m)\). \(F\) es suryectiva: si \(x \in X\), existe \(n\) con \(x \in X_n\) y existe \(m\) con \(x = f_n(m) = F(n, m)\). Como \(\N \times \N\) es numerable, \(X\) es a lo sumo numerable por el teorema anterior; y como contiene al conjunto infinito \(X_1\), es numerable.
\end{proof}

\begin{eg}
  El conjunto \(\Q[x]\) de polinomios con coeficientes racionales es numerable.
\end{eg}

\begin{proof}
  Para cada \(n\), sea \(\Q_n[x] = \{ p \in \Q[x] : \operatorname{gr}(p) \leq n \}\). La función \(\Q^{n+1} \to \Q_n[x]\), \((a_0, \dots, a_n) \mapsto a_0 + a_1 x + \cdots + a_n x^n\), es biyectiva. Como \(\Q^{n+1}\) es numerable (producto finito de numerables) y \(\Q[x] = \bigcup_{n \geq 0} \Q_n[x]\) es unión numerable de conjuntos numerables, \(\Q[x]\) es numerable.
\end{proof}

\begin{definition}[Números algebraicos y trascendentes]
  Un número real es \emph{algebraico} si es raíz de algún polinomio no nulo con coeficientes enteros; por ejemplo, \(\sqrt{2}\) es raíz de \(x^2 - 2\). Un número real que no es algebraico se llama \emph{trascendente}.
\end{definition}

\begin{ejercicio}
  Probar que el conjunto de números algebraicos es numerable. Concluir que existen números trascendentes (de hecho, el conjunto de números trascendentes no es numerable).
\end{ejercicio}
